% ============================================================================
%  SECCIÓN 3.10: CONFIGURACIÓN DEL ENTORNO DE DESARROLLO
% ============================================================================

\section{Configuraci\'on del entorno de desarrollo}

El entorno de desarrollo es el conjunto de herramientas y configuraciones
necesarias para crear, ejecutar y depurar la aplicaci\'on m\'ovil. Para el
presente proyecto se configur\'o un entorno basado en el sistema operativo
Linux, utilizando como editor de c\'odigo y gestor del proyecto la herramienta
Android Studio, que ofrece un entorno de desarrollo integrado (IDE) oficial
para aplicaciones Android y soporte para el desarrollo multiplataforma con
React Native \citep{androidstudio}.

La gesti\'on de paquetes y dependencias se realiza con \texttt{pnpm}, un
gestor de paquetes r\'apido y eficiente que instala las dependencias de forma
determinista mediante el archivo \texttt{pnpm-lock.yaml} \citep{nodejs}. El
dispositivo f\'isico se conecta al equipo mediante depuraci\'on USB y se
habilita la opci\'on de \'origines desconocidos, lo que permite instalar y
ejecutar la aplicaci\'on de forma directa durante las pruebas.

En la tabla se detallan los principales elementos del entorno de desarrollo
configurado en el equipo de trabajo.

\begin{table}[H]
  \centering
  \caption{Elementos del entorno de desarrollo}
  \contenidoConFuente{%
    \begin{tablaAPA}
      \begin{tabular}{@{}p{4.2cm}p{4.6cm}p{7.0cm}@{}}
        \toprule
        \textbf{Componente} & \textbf{Herramienta} & \textbf{Prop\'osito} \\
        \midrule
        Sistema operativo & Linux & Plataforma base del entorno de desarrollo \\
        \addlinespace
        IDE & Android Studio & Edici\'on, gesti\'on y depuraci\'on del proyecto \\
        \addlinespace
        Gestor de paquetes & pnpm & Instalaci\'on y gesti\'on de las
        dependencias del proyecto \\
        \addlinespace
        Runtime de ejecuci\'on & Node.js & Ejecuci\'on de Expo CLI y del
        c\'odigo JavaScript \\
        \addlinespace
        Marco de desarrollo & Expo (React Native) & Creaci\'on y ejecuci\'on de la
        aplicaci\'on multiplataforma \citep{expo} \\
        \addlinespace
        Control de versiones & Git & Registro y trazabilidad de los cambios del
        c\'odigo \citep{gitdocs} \\
        \addlinespace
        Dispositivo de pruebas & Tel\'efono f\'isico (Android) & Ejecuci\'on de la
        aplicaci\'on en un dispositivo real \\
        \bottomrule
      \end{tabular}
    \end{tablaAPA}%
  }{Elaboraci\'on propia en base a la revisi\'on documental.}
\end{table}

Dado que todas las herramientas ya se encuentran instaladas y configuradas en
el equipo de desarrollo, se verific\'o el correcto funcionamiento del entorno
consultando las versiones instaladas de cada componente. En la figura se
presenta la evidencia de la verificaci\'on del entorno de desarrollo.

\begin{figure}[H]
\centering
  \caption{Verificaci\'on del entorno de desarrollo en la terminal}
  \label{fig:entorno-desarrollo}
  \vspace{6pt}
  \begin{lstlisting}[style=terminal]
$ node --version
v24.16.0
$ pnpm --version
11.5.1
$ git --version
git version 2.47.3
$ adb devices
List of devices attached
R58M75TGGVA    device
  \end{lstlisting}
  \par\raggedright\fuenteAPA{Elaboraci\'on propia en base a la revisi\'on documental.}\end{figure}

La salida del comando \texttt{adb devices} confirma que el tel\'efono f\'isico
se encuentra conectado y reconocido por el entorno, listo para la instalaci\'on
y la ejecuci\'on de la aplicaci\'on durante las pruebas.
